
\documentclass{article}
\usepackage{colortbl}
\usepackage{makecell}
\usepackage{multirow}
\usepackage{supertabular}

\begin{document}

\newcounter{utterance}

\centering \large Interaction Transcript for game `dond', experiment `coop\_it', episode 19 with llama{-}3.3{-}70b{-}instruct{-}t1.0.
\vspace{24pt}

{ \footnotesize  \setcounter{utterance}{1}
\setlength{\tabcolsep}{0pt}
\begin{supertabular}{c@{$\;$}|p{.15\linewidth}@{}p{.15\linewidth}p{.15\linewidth}p{.15\linewidth}p{.15\linewidth}p{.15\linewidth}}
    \# & $\;$A & \multicolumn{4}{c}{Game Master} & $\;\:$B\\
    \hline

    \theutterance \stepcounter{utterance}  
    & & \multicolumn{4}{p{0.6\linewidth}}{
        \cellcolor[rgb]{0.9,0.9,0.9}{
            \makecell[{{p{\linewidth}}}]{
                \texttt{\tiny{[P1$\langle$GM]}}
                \texttt{State giocando a un gioco di negoziazione cooperativa in cui dovete accordarvi su come dividere una serie di oggetti tra voi e un altro giocatore.} \\
\\ 
\texttt{Regole:} \\
\texttt{(a) A Lei e all'altro giocatore viene dato un insieme di oggetti. Ognuno di voi riceve anche una funzione di valore segreta, che rappresenta il valore di ciascun tipo di oggetto.} \\
\texttt{(b) Si scambiano messaggi con l'altro giocatore per concordare chi si aggiudica gli oggetti. Potete inviare un massimo di 5 messaggi ciascuno, oppure terminare in anticipo facendo una proposta segreta in qualsiasi momento.} \\
\texttt{(c) A ciascuno di voi viene chiesto di inviare una proposta segreta indicando gli oggetti che desiderate, formattata tra parentesi quadre come segue: "{[}Proposta: <numero> <nome oggetto>, <numero> <nome oggetto>, <...>{]}".} \\
\texttt{(d) Se le vostre proposte sono complementari, cioè ci sono abbastanza oggetti per soddisfare entrambe le proposte, a ciascun giocatore viene assegnato un punteggio basato sulla somma dei valori degli oggetti ricevuti. In caso contrario, entrambi ricevono zero punti.} \\
\texttt{(e) L'obiettivo è massimizzare la somma del proprio punteggio e di quello dell'altro giocatore.} \\
\\ 
\texttt{Cominciamo.} \\
\\ 
\texttt{L'insieme degli oggetti disponibili è:} \\
\\ 
\texttt{2 ganci, 1 fiore, 1 scheda, 1 gomma, 1 spugna} \\
\\ 
\texttt{La funzione valore segreta è:} \\
\\ 
\texttt{gancio: 3, fiore: 3, scheda: 0, gomma: 1, spugna: 0} \\
\\ 
\texttt{IMPORTANTE: i vostri messaggi, a meno che non si tratti di una proposta segreta, vengono trasmessi direttamente all'altro giocatore, quindi non includete alcuna risposta alle regole o testo di annuncio del vostro messaggio. Per fare una proposta segreta, utilizzate il formato indicato. Non utilizzare delle parentesi quadre quando si comunica all'altro giocatore, altrimenti verrà interpretata come una proposta segreta.} \\
            }
        }
    }
    & \\ \\

    \theutterance \stepcounter{utterance}  
    & \multicolumn{4}{p{0.6\linewidth}}{
        \cellcolor[rgb]{0.9,0.9,0.9}{
            \makecell[{{p{\linewidth}}}]{
                \texttt{\tiny{[P1$\rangle$GM]}}
                \texttt{Ciao, vorrei iniziare proponendo una possibile divisione degli oggetti. Ho notato che ci sono alcuni oggetti che potrebbero essere più preziosi per uno di noi due. Che ne pensi di scambiare alcune informazioni sui nostri interessi per gli oggetti disponibili?} \\
            }
        }
    }
    & & \\ \\

    \theutterance \stepcounter{utterance}  
    & & & \multicolumn{2}{p{0.3\linewidth}}{
        \cellcolor[rgb]{0.9,0.9,0.9}{
            \makecell[{{p{\linewidth}}}]{
                \texttt{\tiny{[GM$|$GM]}}
                \texttt{continue} \\
            }
        }
    }
    & & \\ \\

    \theutterance \stepcounter{utterance}  
    & & & \multicolumn{2}{p{0.3\linewidth}}{
        \cellcolor[rgb]{0.9,0.9,0.9}{
            \makecell[{{p{\linewidth}}}]{
                \texttt{\tiny{[GM$|$GM]}}
                \texttt{Ciao, vorrei iniziare proponendo una possibile divisione degli oggetti. Ho notato che ci sono alcuni oggetti che potrebbero essere più preziosi per uno di noi due. Che ne pensi di scambiare alcune informazioni sui nostri interessi per gli oggetti disponibili?} \\
            }
        }
    }
    & & \\ \\

    \theutterance \stepcounter{utterance}  
    & & \multicolumn{4}{p{0.6\linewidth}}{
        \cellcolor[rgb]{0.9,0.9,0.9}{
            \makecell[{{p{\linewidth}}}]{
                \texttt{\tiny{[GM$\rangle$P2]}}
                \texttt{State giocando a un gioco di negoziazione cooperativa in cui dovete accordarvi su come dividere una serie di oggetti tra voi e un altro giocatore.} \\
\\ 
\texttt{Regole:} \\
\texttt{(a) A Lei e all'altro giocatore viene dato un insieme di oggetti. Ognuno di voi riceve anche una funzione di valore segreta, che rappresenta il valore di ciascun tipo di oggetto.} \\
\texttt{(b) Si scambiano messaggi con l'altro giocatore per concordare chi si aggiudica gli oggetti. Potete inviare un massimo di 5 messaggi ciascuno, oppure terminare in anticipo facendo una proposta segreta in qualsiasi momento.} \\
\texttt{(c) A ciascuno di voi viene chiesto di inviare una proposta segreta indicando gli oggetti che desiderate, formattata tra parentesi quadre come segue: "{[}Proposta: <numero> <nome oggetto>, <numero> <nome oggetto>, <...>{]}".} \\
\texttt{(d) Se le vostre proposte sono complementari, cioè ci sono abbastanza oggetti per soddisfare entrambe le proposte, a ciascun giocatore viene assegnato un punteggio basato sulla somma dei valori degli oggetti ricevuti. In caso contrario, entrambi ricevono zero punti.} \\
\texttt{(e) L'obiettivo è massimizzare la somma del proprio punteggio e di quello dell'altro giocatore.} \\
\\ 
\texttt{Cominciamo.} \\
\\ 
\texttt{L'insieme degli oggetti disponibili è:} \\
\\ 
\texttt{2 ganci, 1 fiore, 1 scheda, 1 gomma, 1 spugna} \\
\\ 
\texttt{La funzione valore segreta è:} \\
\\ 
\texttt{gancio: 0, fiore: 2, scheda: 1, gomma: 4, spugna: 3} \\
\\ 
\texttt{IMPORTANTE: i vostri messaggi, a meno che non si tratti di una proposta segreta, vengono trasmessi direttamente all'altro giocatore, quindi non includete alcuna risposta alle regole o testo di annuncio del vostro messaggio. Per fare una proposta segreta, utilizzate il formato indicato. Non utilizzare delle parentesi quadre quando si comunica all'altro giocatore, altrimenti verrà interpretata come una proposta segreta.} \\
\\ 
\\ 
\texttt{Ciao, vorrei iniziare proponendo una possibile divisione degli oggetti. Ho notato che ci sono alcuni oggetti che potrebbero essere più preziosi per uno di noi due. Che ne pensi di scambiare alcune informazioni sui nostri interessi per gli oggetti disponibili?} \\
            }
        }
    }
    & \\ \\

    \theutterance \stepcounter{utterance}  
    & & & \multicolumn{4}{p{0.6\linewidth}}{
        \cellcolor[rgb]{0.9,0.9,0.9}{
            \makecell[{{p{\linewidth}}}]{
                \texttt{\tiny{[GM$\langle$P2]}}
                \texttt{Sembra un'idea interessante. Sono disposto a condividere alcune informazioni sui miei interessi, ma potremmo anche iniziare a discutere gli oggetti specifici che ci interessano di più. Ad esempio, ci sono oggetti che ti sembrano particolarmente utili o interessanti? Forse potremmo iniziare con quelli e vedere se possiamo trovare un accordo.} \\
            }
        }
    }
     \\ \\

    \theutterance \stepcounter{utterance}  
    & & & \multicolumn{2}{p{0.3\linewidth}}{
        \cellcolor[rgb]{0.9,0.9,0.9}{
            \makecell[{{p{\linewidth}}}]{
                \texttt{\tiny{[GM$|$GM]}}
                \texttt{continue} \\
            }
        }
    }
    & & \\ \\

    \theutterance \stepcounter{utterance}  
    & & & \multicolumn{2}{p{0.3\linewidth}}{
        \cellcolor[rgb]{0.9,0.9,0.9}{
            \makecell[{{p{\linewidth}}}]{
                \texttt{\tiny{[GM$|$GM]}}
                \texttt{Sembra un'idea interessante. Sono disposto a condividere alcune informazioni sui miei interessi, ma potremmo anche iniziare a discutere gli oggetti specifici che ci interessano di più. Ad esempio, ci sono oggetti che ti sembrano particolarmente utili o interessanti? Forse potremmo iniziare con quelli e vedere se possiamo trovare un accordo.} \\
            }
        }
    }
    & & \\ \\

    \theutterance \stepcounter{utterance}  
    & & \multicolumn{4}{p{0.6\linewidth}}{
        \cellcolor[rgb]{0.9,0.9,0.9}{
            \makecell[{{p{\linewidth}}}]{
                \texttt{\tiny{[P1$\langle$GM]}}
                \texttt{Sembra un'idea interessante. Sono disposto a condividere alcune informazioni sui miei interessi, ma potremmo anche iniziare a discutere gli oggetti specifici che ci interessano di più. Ad esempio, ci sono oggetti che ti sembrano particolarmente utili o interessanti? Forse potremmo iniziare con quelli e vedere se possiamo trovare un accordo.} \\
            }
        }
    }
    & \\ \\

    \theutterance \stepcounter{utterance}  
    & \multicolumn{4}{p{0.6\linewidth}}{
        \cellcolor[rgb]{0.9,0.9,0.9}{
            \makecell[{{p{\linewidth}}}]{
                \texttt{\tiny{[P1$\rangle$GM]}}
                \texttt{Sì, credo che iniziare con gli oggetti specifici sia un buon approccio. Personalmente, mi sembra che i ganci e il fiore siano piuttosto interessanti. Il fatto che ci siano due ganci potrebbe essere un vantaggio per entrambi, a seconda di come decidiamo di dividerli. Che ne pensi del fiore? Ti interessa molto o potresti essere disposto a lasciarlo a me in cambio di qualcos'altro?} \\
            }
        }
    }
    & & \\ \\

    \theutterance \stepcounter{utterance}  
    & & & \multicolumn{2}{p{0.3\linewidth}}{
        \cellcolor[rgb]{0.9,0.9,0.9}{
            \makecell[{{p{\linewidth}}}]{
                \texttt{\tiny{[GM$|$GM]}}
                \texttt{continue} \\
            }
        }
    }
    & & \\ \\

    \theutterance \stepcounter{utterance}  
    & & & \multicolumn{2}{p{0.3\linewidth}}{
        \cellcolor[rgb]{0.9,0.9,0.9}{
            \makecell[{{p{\linewidth}}}]{
                \texttt{\tiny{[GM$|$GM]}}
                \texttt{Sì, credo che iniziare con gli oggetti specifici sia un buon approccio. Personalmente, mi sembra che i ganci e il fiore siano piuttosto interessanti. Il fatto che ci siano due ganci potrebbe essere un vantaggio per entrambi, a seconda di come decidiamo di dividerli. Che ne pensi del fiore? Ti interessa molto o potresti essere disposto a lasciarlo a me in cambio di qualcos'altro?} \\
            }
        }
    }
    & & \\ \\

    \theutterance \stepcounter{utterance}  
    & & \multicolumn{4}{p{0.6\linewidth}}{
        \cellcolor[rgb]{0.9,0.9,0.9}{
            \makecell[{{p{\linewidth}}}]{
                \texttt{\tiny{[GM$\rangle$P2]}}
                \texttt{Sì, credo che iniziare con gli oggetti specifici sia un buon approccio. Personalmente, mi sembra che i ganci e il fiore siano piuttosto interessanti. Il fatto che ci siano due ganci potrebbe essere un vantaggio per entrambi, a seconda di come decidiamo di dividerli. Che ne pensi del fiore? Ti interessa molto o potresti essere disposto a lasciarlo a me in cambio di qualcos'altro?} \\
            }
        }
    }
    & \\ \\

    \theutterance \stepcounter{utterance}  
    & & & \multicolumn{4}{p{0.6\linewidth}}{
        \cellcolor[rgb]{0.9,0.9,0.9}{
            \makecell[{{p{\linewidth}}}]{
                \texttt{\tiny{[GM$\langle$P2]}}
                \texttt{Il fiore mi sembra un oggetto unico e potenzialmente prezioso. Non sono ancora sicuro di quanto valga per me, ma sono disposto a considerare la possibilità di lasciarlo a te se ottengo qualcosa che mi interessa di più. I ganci, come hai detto, potrebbero essere un buon punto di partenza per una trattativa. Sarei disposto a considerare una divisione dei ganci, ma vorrei sapere cosa pensi della gomma e della spugna. Ti sembrano oggetti utili o non ti interessano particolarmente?} \\
            }
        }
    }
     \\ \\

    \theutterance \stepcounter{utterance}  
    & & & \multicolumn{2}{p{0.3\linewidth}}{
        \cellcolor[rgb]{0.9,0.9,0.9}{
            \makecell[{{p{\linewidth}}}]{
                \texttt{\tiny{[GM$|$GM]}}
                \texttt{continue} \\
            }
        }
    }
    & & \\ \\

    \theutterance \stepcounter{utterance}  
    & & & \multicolumn{2}{p{0.3\linewidth}}{
        \cellcolor[rgb]{0.9,0.9,0.9}{
            \makecell[{{p{\linewidth}}}]{
                \texttt{\tiny{[GM$|$GM]}}
                \texttt{Il fiore mi sembra un oggetto unico e potenzialmente prezioso. Non sono ancora sicuro di quanto valga per me, ma sono disposto a considerare la possibilità di lasciarlo a te se ottengo qualcosa che mi interessa di più. I ganci, come hai detto, potrebbero essere un buon punto di partenza per una trattativa. Sarei disposto a considerare una divisione dei ganci, ma vorrei sapere cosa pensi della gomma e della spugna. Ti sembrano oggetti utili o non ti interessano particolarmente?} \\
            }
        }
    }
    & & \\ \\

    \theutterance \stepcounter{utterance}  
    & & \multicolumn{4}{p{0.6\linewidth}}{
        \cellcolor[rgb]{0.9,0.9,0.9}{
            \makecell[{{p{\linewidth}}}]{
                \texttt{\tiny{[P1$\langle$GM]}}
                \texttt{Il fiore mi sembra un oggetto unico e potenzialmente prezioso. Non sono ancora sicuro di quanto valga per me, ma sono disposto a considerare la possibilità di lasciarlo a te se ottengo qualcosa che mi interessa di più. I ganci, come hai detto, potrebbero essere un buon punto di partenza per una trattativa. Sarei disposto a considerare una divisione dei ganci, ma vorrei sapere cosa pensi della gomma e della spugna. Ti sembrano oggetti utili o non ti interessano particolarmente?} \\
            }
        }
    }
    & \\ \\

    \theutterance \stepcounter{utterance}  
    & \multicolumn{4}{p{0.6\linewidth}}{
        \cellcolor[rgb]{0.9,0.9,0.9}{
            \makecell[{{p{\linewidth}}}]{
                \texttt{\tiny{[P1$\rangle$GM]}}
                \texttt{La gomma e la spugna non mi sembrano oggetti di grande valore per me, ma potrebbero essere utili per te. La gomma in particolare potrebbe essere interessante se hai un uso specifico in mente. La spugna, invece, non mi sembra particolarmente rilevante, mapotrebbe essere utile in determinati contesti. Se fossi disposto a lasciarmi entrambi i ganci, potrei considerare di offrirti la gomma e forse anche la spugna, a patto che tu sia disposto a lasciarmi il fiore. Che ne pensi di questa possibile divisione?} \\
            }
        }
    }
    & & \\ \\

    \theutterance \stepcounter{utterance}  
    & & & \multicolumn{2}{p{0.3\linewidth}}{
        \cellcolor[rgb]{0.9,0.9,0.9}{
            \makecell[{{p{\linewidth}}}]{
                \texttt{\tiny{[GM$|$GM]}}
                \texttt{continue} \\
            }
        }
    }
    & & \\ \\

    \theutterance \stepcounter{utterance}  
    & & & \multicolumn{2}{p{0.3\linewidth}}{
        \cellcolor[rgb]{0.9,0.9,0.9}{
            \makecell[{{p{\linewidth}}}]{
                \texttt{\tiny{[GM$|$GM]}}
                \texttt{La gomma e la spugna non mi sembrano oggetti di grande valore per me, ma potrebbero essere utili per te. La gomma in particolare potrebbe essere interessante se hai un uso specifico in mente. La spugna, invece, non mi sembra particolarmente rilevante, mapotrebbe essere utile in determinati contesti. Se fossi disposto a lasciarmi entrambi i ganci, potrei considerare di offrirti la gomma e forse anche la spugna, a patto che tu sia disposto a lasciarmi il fiore. Che ne pensi di questa possibile divisione?} \\
            }
        }
    }
    & & \\ \\

    \theutterance \stepcounter{utterance}  
    & & \multicolumn{4}{p{0.6\linewidth}}{
        \cellcolor[rgb]{0.9,0.9,0.9}{
            \makecell[{{p{\linewidth}}}]{
                \texttt{\tiny{[GM$\rangle$P2]}}
                \texttt{La gomma e la spugna non mi sembrano oggetti di grande valore per me, ma potrebbero essere utili per te. La gomma in particolare potrebbe essere interessante se hai un uso specifico in mente. La spugna, invece, non mi sembra particolarmente rilevante, mapotrebbe essere utile in determinati contesti. Se fossi disposto a lasciarmi entrambi i ganci, potrei considerare di offrirti la gomma e forse anche la spugna, a patto che tu sia disposto a lasciarmi il fiore. Che ne pensi di questa possibile divisione?} \\
            }
        }
    }
    & \\ \\

    \theutterance \stepcounter{utterance}  
    & & & \multicolumn{4}{p{0.6\linewidth}}{
        \cellcolor[rgb]{0.9,0.9,0.9}{
            \makecell[{{p{\linewidth}}}]{
                \texttt{\tiny{[GM$\langle$P2]}}
                \texttt{L'idea di lasciarti entrambi i ganci in cambio della gomma e forse della spugna, oltre a tenere il fiore per te, sembra un buon compromesso. La gomma potrebbe essere particolarmente utile per me, e anche la spugna potrebbe essere interessante se riuscissi a trovare un uso appropriato. Tuttavia, vorrei assicurarmi di non stare cedendo troppo. La scheda mi sembra un oggetto abbastanza interessante, non hai detto molto al riguardo. Ti interessa la scheda o potresti essere disposto a includerla nell'accordo?} \\
            }
        }
    }
     \\ \\

    \theutterance \stepcounter{utterance}  
    & & & \multicolumn{2}{p{0.3\linewidth}}{
        \cellcolor[rgb]{0.9,0.9,0.9}{
            \makecell[{{p{\linewidth}}}]{
                \texttt{\tiny{[GM$|$GM]}}
                \texttt{continue} \\
            }
        }
    }
    & & \\ \\

    \theutterance \stepcounter{utterance}  
    & & & \multicolumn{2}{p{0.3\linewidth}}{
        \cellcolor[rgb]{0.9,0.9,0.9}{
            \makecell[{{p{\linewidth}}}]{
                \texttt{\tiny{[GM$|$GM]}}
                \texttt{L'idea di lasciarti entrambi i ganci in cambio della gomma e forse della spugna, oltre a tenere il fiore per te, sembra un buon compromesso. La gomma potrebbe essere particolarmente utile per me, e anche la spugna potrebbe essere interessante se riuscissi a trovare un uso appropriato. Tuttavia, vorrei assicurarmi di non stare cedendo troppo. La scheda mi sembra un oggetto abbastanza interessante, non hai detto molto al riguardo. Ti interessa la scheda o potresti essere disposto a includerla nell'accordo?} \\
            }
        }
    }
    & & \\ \\

    \theutterance \stepcounter{utterance}  
    & & \multicolumn{4}{p{0.6\linewidth}}{
        \cellcolor[rgb]{0.9,0.9,0.9}{
            \makecell[{{p{\linewidth}}}]{
                \texttt{\tiny{[P1$\langle$GM]}}
                \texttt{L'idea di lasciarti entrambi i ganci in cambio della gomma e forse della spugna, oltre a tenere il fiore per te, sembra un buon compromesso. La gomma potrebbe essere particolarmente utile per me, e anche la spugna potrebbe essere interessante se riuscissi a trovare un uso appropriato. Tuttavia, vorrei assicurarmi di non stare cedendo troppo. La scheda mi sembra un oggetto abbastanza interessante, non hai detto molto al riguardo. Ti interessa la scheda o potresti essere disposto a includerla nell'accordo?} \\
            }
        }
    }
    & \\ \\

    \theutterance \stepcounter{utterance}  
    & \multicolumn{4}{p{0.6\linewidth}}{
        \cellcolor[rgb]{0.9,0.9,0.9}{
            \makecell[{{p{\linewidth}}}]{
                \texttt{\tiny{[P1$\rangle$GM]}}
                \texttt{La scheda non mi interessa particolarmente, quindi sarei più che disposto a includerla nell'accordo. In realtà, non la considero molto importante, quindi se lasciarti la scheda insieme alla gomma e alla spugna in cambio dei ganci e del permesso di tenere il fiore, mi sembra un buon equilibrio. Sembra che stiamo arrivando a un accordo, quindi potremmo iniziare a pensare a come formalizzare la nostra proposta. Che ne pensi di fare una proposta segreta che rifletta questo accordo?} \\
            }
        }
    }
    & & \\ \\

    \theutterance \stepcounter{utterance}  
    & & & \multicolumn{2}{p{0.3\linewidth}}{
        \cellcolor[rgb]{0.9,0.9,0.9}{
            \makecell[{{p{\linewidth}}}]{
                \texttt{\tiny{[GM$|$GM]}}
                \texttt{continue} \\
            }
        }
    }
    & & \\ \\

    \theutterance \stepcounter{utterance}  
    & & & \multicolumn{2}{p{0.3\linewidth}}{
        \cellcolor[rgb]{0.9,0.9,0.9}{
            \makecell[{{p{\linewidth}}}]{
                \texttt{\tiny{[GM$|$GM]}}
                \texttt{La scheda non mi interessa particolarmente, quindi sarei più che disposto a includerla nell'accordo. In realtà, non la considero molto importante, quindi se lasciarti la scheda insieme alla gomma e alla spugna in cambio dei ganci e del permesso di tenere il fiore, mi sembra un buon equilibrio. Sembra che stiamo arrivando a un accordo, quindi potremmo iniziare a pensare a come formalizzare la nostra proposta. Che ne pensi di fare una proposta segreta che rifletta questo accordo?} \\
            }
        }
    }
    & & \\ \\

    \theutterance \stepcounter{utterance}  
    & & \multicolumn{4}{p{0.6\linewidth}}{
        \cellcolor[rgb]{0.9,0.9,0.9}{
            \makecell[{{p{\linewidth}}}]{
                \texttt{\tiny{[GM$\rangle$P2]}}
                \texttt{La scheda non mi interessa particolarmente, quindi sarei più che disposto a includerla nell'accordo. In realtà, non la considero molto importante, quindi se lasciarti la scheda insieme alla gomma e alla spugna in cambio dei ganci e del permesso di tenere il fiore, mi sembra un buon equilibrio. Sembra che stiamo arrivando a un accordo, quindi potremmo iniziare a pensare a come formalizzare la nostra proposta. Che ne pensi di fare una proposta segreta che rifletta questo accordo?} \\
            }
        }
    }
    & \\ \\

    \theutterance \stepcounter{utterance}  
    & & & \multicolumn{4}{p{0.6\linewidth}}{
        \cellcolor[rgb]{0.9,0.9,0.9}{
            \makecell[{{p{\linewidth}}}]{
                \texttt{\tiny{[GM$\langle$P2]}}
                \texttt{Sì, credo che abbiamo trovato un buon equilibrio. Sono d'accordo con la divisione che hai descritto: tu prendi i ganci e il fiore, e io prendo la gomma, la spugna e la scheda. È arrivato il momento di formalizzare la nostra proposta. Preparerò la mia proposta segreta secondo l'accordo che abbiamo discusso. {[}Proposta: 1 gomma, 1 spugna, 1 scheda{]}} \\
            }
        }
    }
     \\ \\

    \theutterance \stepcounter{utterance}  
    & & & \multicolumn{2}{p{0.3\linewidth}}{
        \cellcolor[rgb]{0.9,0.9,0.9}{
            \makecell[{{p{\linewidth}}}]{
                \texttt{\tiny{[GM$|$GM]}}
                \texttt{proposal} \\
            }
        }
    }
    & & \\ \\

    \theutterance \stepcounter{utterance}  
    & & & \multicolumn{2}{p{0.3\linewidth}}{
        \cellcolor[rgb]{0.9,0.9,0.9}{
            \makecell[{{p{\linewidth}}}]{
                \texttt{\tiny{[GM$|$GM]}}
                [0, 0, 1, 1, 1]
            }
        }
    }
    & & \\ \\

    \theutterance \stepcounter{utterance}  
    & & \multicolumn{4}{p{0.6\linewidth}}{
        \cellcolor[rgb]{0.9,0.9,0.9}{
            \makecell[{{p{\linewidth}}}]{
                \texttt{\tiny{[P1$\langle$GM]}}
                \texttt{L'altro giocatore ha presentato una proposta segreta. Ora tocca a Lei presentare una sua proposta.} \\
            }
        }
    }
    & \\ \\

    \theutterance \stepcounter{utterance}  
    & \multicolumn{4}{p{0.6\linewidth}}{
        \cellcolor[rgb]{0.9,0.9,0.9}{
            \makecell[{{p{\linewidth}}}]{
                \texttt{\tiny{[P1$\rangle$GM]}}
                \texttt{{[}Proposta: 2 ganci, 1 fiore, 0 scheda, 0 gomma, 0 spugna{]}} \\
            }
        }
    }
    & & \\ \\

    \theutterance \stepcounter{utterance}  
    & & & \multicolumn{2}{p{0.3\linewidth}}{
        \cellcolor[rgb]{0.9,0.9,0.9}{
            \makecell[{{p{\linewidth}}}]{
                \texttt{\tiny{[GM$|$GM]}}
                \texttt{proposal} \\
            }
        }
    }
    & & \\ \\

    \theutterance \stepcounter{utterance}  
    & & & \multicolumn{2}{p{0.3\linewidth}}{
        \cellcolor[rgb]{0.9,0.9,0.9}{
            \makecell[{{p{\linewidth}}}]{
                \texttt{\tiny{[GM$|$GM]}}
                [2, 1, 0, 0, 0]
            }
        }
    }
    & & \\ \\

    \theutterance \stepcounter{utterance}  
    & & & \multicolumn{2}{p{0.3\linewidth}}{
        \cellcolor[rgb]{0.9,0.9,0.9}{
            \makecell[{{p{\linewidth}}}]{
                \texttt{\tiny{[GM$|$GM]}}
                [[2, 1, 0, 0, 0], [0, 0, 1, 1, 1]]
            }
        }
    }
    & & \\ \\

\end{supertabular}
}

\end{document}
